% Options for packages loaded elsewhere
\PassOptionsToPackage{unicode}{hyperref}
\PassOptionsToPackage{hyphens}{url}
\PassOptionsToPackage{space}{xeCJK}
\documentclass[
  a4paper,11pt]{ltjsarticle}
\usepackage{xcolor}
\usepackage[margin=1in]{geometry}
\usepackage{amsmath,amssymb}
\usepackage{graphicx}
\setcounter{secnumdepth}{-\maxdimen} % remove section numbering
\usepackage{iftex}
\ifPDFTeX
  \usepackage[T1]{fontenc}
  \usepackage[utf8]{inputenc}
  \usepackage{textcomp}
\else
  \usepackage{unicode-math}
  \defaultfontfeatures{Scale=MatchLowercase}
  \defaultfontfeatures[\rmfamily]{Ligatures=TeX,Scale=1}
\fi
\usepackage{lmodern}
\ifPDFTeX\else
  \setmainfont[]{Times New Roman}
  \ifXeTeX
    \usepackage{xeCJK}
    \setCJKmainfont[]{Hiragino Mincho ProN}
  \fi
  \ifLuaTeX
    \usepackage[]{luatexja-fontspec}
    \setmainjfont[]{Hiragino Mincho ProN}
  \fi
\fi
\IfFileExists{upquote.sty}{\usepackage{upquote}}{}
\IfFileExists{microtype.sty}{%
  \usepackage[]{microtype}
  \UseMicrotypeSet[protrusion]{basicmath}
}{}
\makeatletter
\@ifundefined{KOMAClassName}{%
  \IfFileExists{parskip.sty}{%
    \usepackage{parskip}
  }{%
    \setlength{\parindent}{0pt}
    \setlength{\parskip}{6pt plus 2pt minus 1pt}}
}{% if KOMA class
  \KOMAoptions{parskip=half}}
\makeatother
\usepackage{longtable,booktabs,array}
\newcounter{none}
\usepackage{calc}
\usepackage{etoolbox}
\makeatletter
\patchcmd\longtable{\par}{\if@noskipsec\mbox{}\fi\par}{}{}
\makeatother
\IfFileExists{footnotehyper.sty}{\usepackage{footnotehyper}}{\usepackage{footnote}}
\makesavenoteenv{longtable}
\usepackage{bookmark}
\IfFileExists{xurl.sty}{\usepackage{xurl}}{}
\urlstyle{same}
\usepackage{hyperref}
\hypersetup{
  hidelinks,
  pdfcreator={LaTeX via pandoc}}
\setlength{\emergencystretch}{3em}
\providecommand{\tightlist}{%
  \setlength{\itemsep}{0pt}\setlength{\parskip}{0pt}}
\usepackage{amsthm}
\newtheorem{theorem}{定理}[section]
\newtheorem{proposition}[theorem]{命題}
\newtheorem{lemma}[theorem]{補題}
\newtheorem{corollary}[theorem]{系}
\theoremstyle{definition}
\newtheorem{definition}[theorem]{定義}
\theoremstyle{remark}
\newtheorem{remark}[theorem]{注意}
\newtheorem{observation}[theorem]{観察}
\begin{document}
\section{4次元整数格子上の形不変移動波------存在条件・保存構造・自己整合ループ}\label{ux6b21ux5143ux6574ux6570ux683cux5b50ux4e0aux306eux5f62ux4e0dux5909ux79fbux52d5ux6ce2ux5b58ux5728ux6761ux4ef6ux4fddux5b58ux69cbux9020ux81eaux5df1ux6574ux5408ux30ebux30fcux30d7}

\textbf{著者}: 木原 範昭 (Noriaki Kihara)\\
\textbf{所属}: WF System Co., Ltd.\\
\textbf{ORCID}: 0009-0004-6753-4020\\
\textbf{版}: v1.0\\
\textbf{日付}: 2026 年 4 月

\begin{center}\rule{0.5\linewidth}{0.5pt}\end{center}

\subsection{要旨}\label{ux8981ux65e8}

\(n\) 次元整数格子 \(\mathbb{Z}^n\) 上で、平行移動に対して振幅分布
\(|\psi|\)
を厳密に保存する波------形不変移動波------が整合的に存在しうる条件を、球面整合性・乗法閉性・モード分解可能性の3軸で整理する。球面定常波の存在条件
(D1)--(D4) は平方数次元 \(n \in \{4, 9, 16, \ldots\}\)
で成立するが、結合的乗法構造（Frobenius
の定理）と明示的モード計数（Jacobi
四平方公式）を同時に要求すると、\(n = 4\) が \textbf{唯一の次元}
として特定される。この格子上で定義される形不変移動波は、ラベル付き対象としての個別性・共存性・交差不変性・合成則・対称群分類を同時に満たし、接束
\(\mathbb{Z}^4 \times \mathbb{Z}^4\)
への拡張により離散流を内蔵し、直交相互作用波による自己帰還ループが
Nyquist
安定性から速度上限を導出する。さらに、親空間からの中心投影が線形系に非線形効果を時間概念なしで導入する唯一の幾何学的メカニズムであることを示し、これらの構造が
\textbf{de Sitter 時空 \(\mathrm{dS}_4\)}
を幾何学的必然として指し示すことを論じる。

\textbf{キーワード}: 形不変移動波, 整数格子, Hurwitz 四元数環, Jacobi
四平方公式, Nyquist 安定性, 中心投影, de Sitter 時空

\begin{center}\rule{0.5\linewidth}{0.5pt}\end{center}

\subsection{§1
導入と主張範囲}\label{ux5c0eux5165ux3068ux4e3bux5f35ux7bc4ux56f2}

本論文は、\(n\) 次元整数格子 \(\mathbb{Z}^n\)
上で平行移動に対して波形を厳密に保つ波------形不変移動波------がどのような次元依存性を持つかを、純粋に幾何学的・数論的・代数的な事実として記述する。

本論文の主張は以下の4点に限定される:

\begin{enumerate}
\def\labelenumi{\arabic{enumi}.}
\tightlist
\item
  球面定常波の存在条件 (D1)--(D4) と、形不変移動波の構造条件 (F1)--(F5)
  の交差により、\(n = 4\) が唯一の整合次元として特定される。
\item
  \(\mathbb{Z}^4\)
  上の形不変移動波は、ラベル付き対象として個別性・共存性・交差不変性・合成則・対称群分類の5条件を同時に満たす。
\item
  接束拡張と直交相互作用波により自己帰還ループが構成され、Nyquist
  安定性条件から速度上限が幾何学的に導出される。
\item
  親空間からの中心投影が、線形系に非線形効果を時間概念なしで導入する唯一の幾何学的メカニズムである。
\end{enumerate}

以下の要素は \textbf{本論文の主張範囲外} である:

\begin{itemize}
\tightlist
\item
  力学的安定性（時間発展に対する固有値解析、エネルギー汎関数の凸性）
\item
  非線形波動方程式の解としての厳密な定式化
\item
  物理的粒子との同一視、量子場理論の再現
\end{itemize}

本論文で「形不変移動波」と呼ぶ対象は、力学的含意を持つ「ソリトン」とは明確に区別される。「形不変」とは平行移動に対して
\(|\psi|\)
が保存されるという幾何学的性質のみを指す。「移動」とは平行移動という格子自己同型に対する共変性のみを指し、運動や伝播の動力学的概念は含まない。

\begin{center}\rule{0.5\linewidth}{0.5pt}\end{center}

\subsection{§2
4次元整数格子の数論的特権性}\label{ux6b21ux5143ux6574ux6570ux683cux5b50ux306eux6570ux8ad6ux7684ux7279ux6a29ux6027}

\subsubsection{§2.1 Lagrange
四平方定理と半径²集合の完全性}\label{lagrange-ux56dbux5e73ux65b9ux5b9aux7406ux3068ux534aux5f84uxb2ux96c6ux5408ux306eux5b8cux5168ux6027}

\(\mathbb{Z}^n\) の格子点間の二乗距離の集合
\(\{|\mathbf{x}|^2 : \mathbf{x} \in \mathbb{Z}^n\}\) を考える。

\(n = 2\) では \(3, 6, 7, \ldots\) が欠落し、\(n = 3\) では
\(4^k(8m+7)\) の形の整数が欠落する。Lagrange
の四平方定理により、\(n \geq 4\) で初めて

\[\{|\mathbf{x}|^2 : \mathbf{x} \in \mathbb{Z}^n\} = \mathbb{Z}_{\geq 0}\]

が成立する。\(n = 4\) はこの完全性が成立する \textbf{最小次元} である。

\subsubsection{§2.2 Jacobi
四平方公式と球面格子点の非空性}\label{jacobi-ux56dbux5e73ux65b9ux516cux5f0fux3068ux7403ux9762ux683cux5b50ux70b9ux306eux975eux7a7aux6027}

半径 \(\sqrt{N}\) の超球面上の \(\mathbb{Z}^n\) 格子点数 \(r_n(N)\)
について、Jacobi の四平方公式は

\[r_4(N) = 8 \sum_{\substack{d \mid N \\ 4 \nmid d}} d > 0 \quad (N \geq 1)\]

を与える。\(r_4(N)\) は全ての \(N \geq 1\)
で正であり、しかも約数和関数で \textbf{閉じた明示式} を持つ。

\(r_3(N)\) は \(N = 4^k(8m+7)\) で零点を持ち、\(r_5, r_6, r_7\)
の明示式は保型形式の尖点形式部分を含み遥かに複雑である。「全 \(N\)
で正、かつ約数和で書ける」を両立するのは \(r_4\) のみである。

\subsubsection{§2.3
等成分方向の整数性}\label{ux7b49ux6210ux5206ux65b9ux5411ux306eux6574ux6570ux6027}

\(n\) 次元単位超立方体の主対角線の長さは \(\sqrt{n}\) であり、等成分方向
\((\pm k, \pm k, \ldots, \pm k)\) に沿った格子線分の長さは
\(\sqrt{n} \cdot k\) である。\(\sqrt{n}\) が整数になるのは \(n\)
が平方数のときに限られ、\(n = 4\)（\(\sqrt{4} = 2\)）はその
\textbf{最小の非自明な次元} である。

\subsubsection{§2.4 Hurwitz
四元数環の乗法構造}\label{hurwitz-ux56dbux5143ux6570ux74b0ux306eux4e57ux6cd5ux69cbux9020}

四元数体 \(\mathbb{H}\) の中に

\[\mathcal{H} = \mathbb{Z}^4 \cup \left(\mathbb{Z}^4 + \left(\tfrac{1}{2}, \tfrac{1}{2}, \tfrac{1}{2}, \tfrac{1}{2}\right)\right)\]

を取ると、これはノルム乗法性 \(|xy| = |x| \cdot |y|\) を満たす
\textbf{左ユークリッド整域} となる。

Frobenius の定理より、実数上の有限次元結合的可除代数は
\(\mathbb{R}, \mathbb{C}, \mathbb{H}\) に限られる。乗法的整数格子は
\(\mathbb{Z}, \mathbb{Z}[i], \mathcal{H}\) のみであり、その次元は
\(1, 2, 4\) である。\(n = 4\) は結合的乗法構造を持つ \textbf{最大次元}
である。

\begin{center}\rule{0.5\linewidth}{0.5pt}\end{center}

\subsection{§3 球面定常波の存在条件
(D1)--(D4)}\label{ux7403ux9762ux5b9aux5e38ux6ce2ux306eux5b58ux5728ux6761ux4ef6-d1d4}

\(\mathbb{Z}^n\) 上で球面族 \(\{S(\sqrt{N})\}_{N \geq 0}\)
と整合する定常波が存在する条件を以下の4つに分解する。

\textbf{(D1) 節半径の整数性}:
節を成す球面族の半径列が、等成分方向に沿って整数間隔で並ぶ。\(\sqrt{n} \in \mathbb{Z}\)
を要求。

\textbf{(D2) 節球面の非空性}:
任意の節半径に対し、その球面上に格子点が存在する。\(r_n(N) > 0\) for all
\(N \geq 1\) を要求。

\textbf{(D3) 節半径²集合の完全性}: 節半径²の取りうる値が
\(\mathbb{Z}_{\geq 0}\)
全体を覆う。\(\{|\mathbf{x}|^2 : \mathbf{x} \in \mathbb{Z}^n\} = \mathbb{Z}_{\geq 0}\)
を要求。

\textbf{(D4) 波数双対性}:
実空間の節格子と波数空間が同型な離散構造を持つ（\(\mathbb{Z}^n\)
の自己双対性）。

\textbf{次元依存性}:

{\def\LTcaptype{none} % do not increment counter
\begin{longtable}[]{@{}llllll@{}}
\toprule\noalign{}
条件 & \(n=2\) & \(n=3\) & \(n=4\) & \(n=9\) & \(n \geq 5\) 一般 \\
\midrule\noalign{}
\endhead
\bottomrule\noalign{}
\endlastfoot
(D1) & × (\(\sqrt{2}\)) & × (\(\sqrt{3}\)) & ✓ & ✓ & 平方数のみ \\
(D2) & × & × & ✓ & ✓ & \(n \geq 4\) \\
(D3) & × & × & ✓ & ✓ & \(n \geq 4\) \\
(D4) & ✓ & ✓ & ✓ & ✓ & 全 \(n\) \\
\end{longtable}
}

\textbf{命題 3.1}: (D1)--(D4) を同時に満たす次元の集合は
\(\{k^2 : k \in \mathbb{Z}_{\geq 2}\} = \{4, 9, 16, 25, \ldots\}\)
であり、最小元は \(n = 4\) である。

\begin{center}\rule{0.5\linewidth}{0.5pt}\end{center}

\subsection{§4 形不変移動波の定義と構造条件
(F1)--(F5)}\label{ux5f62ux4e0dux5909ux79fbux52d5ux6ce2ux306eux5b9aux7fa9ux3068ux69cbux9020ux6761ux4ef6-f1f5}

\subsubsection{§4.1
形不変移動波の定義}\label{ux5f62ux4e0dux5909ux79fbux52d5ux6ce2ux306eux5b9aux7fa9}

\textbf{定義 4.1}: 格子 \(\Lambda \subset \mathbb{R}^n\) 上の関数
\(\psi : \Lambda \to \mathbb{C}\) が \textbf{形不変移動波}
であるとは、ある波数 \(\mathbf{k} \in \Lambda^*\) と任意の平行移動
\(\mathbf{a} \in \Lambda\) に対し、

\[\psi(\mathbf{x} + \mathbf{a}) = e^{i\langle \mathbf{k}, \mathbf{a} \rangle} \cdot \psi(\mathbf{x}), \quad \langle \mathbf{k}, \mathbf{a} \rangle \in \mathbb{Z}\]

が成立すること。

この定義は時間も方程式もエネルギーも含まず、格子・内積・指数写像のみで記述される。

\subsubsection{§4.2 構造条件}\label{ux69cbux9020ux6761ux4ef6}

形不変移動波が「整合的に存在する」ための幾何学的条件:

\textbf{(F1) 位相整合}:
\(\langle \mathbf{k}, \mathbf{a} \rangle \in \mathbb{Z}\) for all
\(\mathbf{a} \in \Lambda\) \(\Leftrightarrow\)
\(\mathbf{k} \in \Lambda^*\)

\textbf{(F2) 振幅保存}:
\(|\psi(\mathbf{x} + \mathbf{a})| = |\psi(\mathbf{x})|\)（定義の直接的帰結）

\textbf{(F3) 合成閉性}: \(\psi_1 \cdot \psi_2\) が再び形不変移動波
\(\Leftrightarrow\) \(\mathbf{k}_1 + \mathbf{k}_2 \in \Lambda^*\)

\textbf{(F4) 球面整合}: 波数の大きさ \(|\mathbf{k}|\)
が球面族の半径列に乗る \(\Leftrightarrow\) (D1)--(D4)

\textbf{(F5) 乗法閉性}: \(\psi(\mathbf{x}) \cdot \psi(\mathbf{y})\)
を格子上の積 \(\mathbf{x} \cdot \mathbf{y}\) の波として書ける
\(\Leftrightarrow\) \(\Lambda\) 上の結合的乗法構造

\subsubsection{§4.3
条件の次元的篩い分け}\label{ux6761ux4ef6ux306eux6b21ux5143ux7684ux7be9ux3044ux5206ux3051}

\begin{itemize}
\tightlist
\item
  (F1)(F2)(F3): \(\Lambda = \mathbb{Z}^n\), \(\Lambda^* = \mathbb{Z}^n\)
  で全 \(n\) に成立
\item
  (F4): \(n \in \{4, 9, 16, 25, \ldots\}\)（平方数次元）
\item
  (F5): \(n \in \{1, 2, 4\}\)（\(\mathbb{R}, \mathbb{C}, \mathbb{H}\);
  八元数 \(n = 8\) は非結合的で除外）
\end{itemize}

\textbf{定理 4.2}: \((F4) \cap (F5) = \{4\}\).

形不変移動波が球面整合性と乗法閉性を同時に満たす整数格子は
\(\mathbb{Z}^4\) に \textbf{限られる}。

\begin{center}\rule{0.5\linewidth}{0.5pt}\end{center}

\subsection{§5 保存量と不変構造
(C1)--(C8)}\label{ux4fddux5b58ux91cfux3068ux4e0dux5909ux69cbux9020-c1c8}

形不変移動波
\(\psi(\mathbf{x} + \mathbf{a}) = e^{i\langle \mathbf{k}, \mathbf{a} \rangle} \cdot \psi(\mathbf{x})\)
の定義から、平行移動 \(\mathbf{a} \in \mathbb{Z}^4\) に対して以下の量が
\textbf{厳密に保存される}。

\textbf{(C1) 振幅分布}:
\(|\psi(\mathbf{x} + \mathbf{a})| = |\psi(\mathbf{x})|\)

\textbf{(C2) 節集合配置}: \(\{\mathbf{x} : \psi(\mathbf{x}) = 0\}\) が
\(\mathbf{a}\) 平行移動で自己同型

\textbf{(C3) 位相差構造}:
\(\langle \mathbf{k}, \mathbf{x}_i - \mathbf{x}_j \rangle \in \mathbb{Z}\)
が任意の格子点対で保存

\textbf{(C4) 波数の大きさ}: \(|\mathbf{k}|\) が平行移動で不変

\textbf{(C5) モード多重度}: \(r_4(|\mathbf{k}|^2)\) が平行移動で不変

4次元では、平行移動を超えた対称群に対する追加の不変構造が存在する:

\textbf{(C6) 四元数乗法に対する保存}: Hurwitz 環の単位元
\(u\)（\(|u| = 1\)）による \(\mathbf{x} \mapsto \mathbf{x} \cdot u\)
に対し、\(|\psi|\)
の球面平均が保存。\(\mathrm{SO}(4) \cong (\mathrm{SU}(2) \times \mathrm{SU}(2))/\mathbb{Z}_2\)
の作用として記述される。

\textbf{(C7) \(D_4\) triality に対する保存}:
ベクトル・スピノル・共役スピノル表現の巡回置換に対し、テータ級数
\(\theta_{D_4}(\tau)\) が不変。

\textbf{(C8) モジュラー変換に対する保存}:
\(\theta_{\mathbb{Z}^4}(\tau)\) が \(\Gamma_0(4)\)
のもとで重み2のモジュラー形式として変換。

\textbf{命題 5.1}: (C1)--(C5) は全次元で成立するが、(C6)--(C8) は
\(n = 4\)
でのみ成立する。4次元の形不変移動波は、平行移動・四元数乗法・三重双対・モジュラー変換のすべてに対する不変量を同時に保持する
\textbf{唯一の構造} である。

\begin{center}\rule{0.5\linewidth}{0.5pt}\end{center}

\subsection{§6 複数波の共存と交差不変性
(G1)--(G5)}\label{ux8907ux6570ux6ce2ux306eux5171ux5b58ux3068ux4ea4ux5deeux4e0dux5909ux6027-g1g5}

\subsubsection{§6.1
重ね合わせと完全分離}\label{ux91cdux306dux5408ux308fux305bux3068ux5b8cux5168ux5206ux96e2}

\(\mathbb{Z}^4\) 上の複数の形不変移動波
\(\psi_{\mathbf{k}_j}(\mathbf{x}) = e^{i\langle \mathbf{k}_j, \mathbf{x} \rangle}\)
を考える。

重ね合わせ
\(\Psi(\mathbf{x}) = \sum_j c_j \cdot \psi_{\mathbf{k}_j}(\mathbf{x})\)
から各成分を取り出す操作は、\(\mathbb{Z}^4\) 上の離散 Fourier
変換の一意性により厳密に実行可能である:

\[c_j = \lim_{N \to \infty} \frac{1}{(2N+1)^4} \sum_{\mathbf{x} \in [-N,N]^4} \Psi(\mathbf{x}) \cdot e^{-i\langle \mathbf{k}_j, \mathbf{x} \rangle}\]

\(\mathbb{Z}^4\)
の双対基底が完全直交系を成すため、任意の有限・可算個のラベル付き波が
\textbf{完全分離} される。

\subsubsection{§6.2
粒子的条件の幾何学的実現}\label{ux7c92ux5b50ux7684ux6761ux4ef6ux306eux5e7eux4f55ux5b66ux7684ux5b9fux73fe}

形不変移動波 \(\psi_{\mathbf{k}}(\mathbf{x})\) を「ラベル \(\mathbf{k}\)
で識別される幾何学的対象」と見なすと、4次元では以下が成立する:

\textbf{(G1) 個別性}: 各 \(\mathbf{k} \in \mathbb{Z}^4\)
が一意のラベルを与え、対象が可算離散集合で識別される。

\textbf{(G2) 共存性}:
任意個数のラベル付き対象が同一空間に整合的に共存する（重ね合わせの閉性）。

\textbf{(G3) 交差不変性}: 複数対象が同じ領域を占めても、Fourier
分解で各ラベルが \textbf{完全に回復} される。

\textbf{(G4) 合成則}: ラベルの加法
\(\mathbf{k}_1 + \mathbf{k}_2 \in \mathbb{Z}^4\)
により合成が定義され、Hurwitz 環の乗法として閉じる。

\textbf{(G5) 分類体系}: ラベルが \(\mathrm{SO}(4)\), \(D_4\) triality,
モジュラー変換のもとで群論的に整理される。

\textbf{定理 6.1}: (G1)--(G5)
を同時に満たす形で整数格子上に形不変移動波が存在しうる唯一の次元は
\(n = 4\) である。

\textbf{注意}: (G1)--(G5)
は粒子に要求される幾何学的条件の構造的対応を内蔵するが、本論文は「幾何学的構造が条件を満たす」という事実のみを主張し、「したがって粒子である」という同一視は行わない。

\begin{center}\rule{0.5\linewidth}{0.5pt}\end{center}

\subsection{§7 ラベル構造
(L1)--(L5)}\label{ux30e9ux30d9ux30ebux69cbux9020-l1l5}

形不変移動波の波数ベクトル \(\mathbf{k} \in \mathbb{Z}^4\)
は整数値ラベルとして機能する。4次元では、このラベル集合が以下の5層の構造を同時に持つ。

\textbf{(L1) 群構造}: ラベル集合 \(\mathbb{Z}^4\)
が四元数加法群として閉じ、さらに Hurwitz 環 \(\mathcal{H}\)
の乗法群構造を持つ。ラベル間の加算・乗算が群論的に閉じる。

\textbf{(L2) 幾何学的階層}: \(|\mathbf{k}|^2\) が整数値を取り、Jacobi
四平方公式 \(r_4(|\mathbf{k}|^2)\)
によりラベル空間が球面殻ごとに自然に階層化される。各殻 \(S(\sqrt{N})\)
上のラベル数が約数和で明示的に与えられる。

\textbf{(L3) 対称群作用}:
\(\mathrm{SO}(4) \cong (\mathrm{SU}(2) \times \mathrm{SU}(2))/\mathbb{Z}_2\)
がラベル集合に作用し、ラベルが4次元回転群の表現として整理される。

\textbf{(L4) 三重双対 (triality)}: \(D_4\) triality
によりラベルが3種類の表現（ベクトル・スピノル・共役スピノル）に分類され、巡回置換可能。

\textbf{(L5) モジュラー変換}: テータ級数 \(\theta_{\mathbb{Z}^4}(\tau)\)
が \(\Gamma_0(4)\)
のもとでモジュラー変換し、ラベル分布のスケール双対性が成立する。

\textbf{命題 7.1}: (L1)--(L5) のすべてを同時に満たすラベル構造は
\(\mathbb{Z}^4\) に限られる。(L1) は \(n \in \{1, 2, 4\}\)、(L2) は
\(n = 4\) でのみ簡潔な閉形式、(L4) は \(n = 4\) でのみ成立する。

\begin{center}\rule{0.5\linewidth}{0.5pt}\end{center}

\subsection{§8
接束拡張と離散流}\label{ux63a5ux675fux62e1ux5f35ux3068ux96e2ux6563ux6d41}

\subsubsection{§8.1
ラベル空間の接束拡張}\label{ux30e9ux30d9ux30ebux7a7aux9593ux306eux63a5ux675fux62e1ux5f35}

形不変移動波のラベル \(\mathbf{k} \in \mathbb{Z}^4\)
を、ラベルと変化率の組
\((\mathbf{k}, \Delta\mathbf{k}) \in \mathbb{Z}^4 \times \mathbb{Z}^4\)
に拡張する。これはラベル空間 \(\mathbb{Z}^4\) の \textbf{離散接束}
\(T\mathbb{Z}^4 = \mathbb{Z}^4 \times \mathbb{Z}^4\) に対応する。

拡張された波は

\[\psi_{(\mathbf{k}, \Delta\mathbf{k})}(\mathbf{x}, t) = e^{i\langle \mathbf{k} + t \cdot \Delta\mathbf{k},\; \mathbf{x} \rangle}\]

と書ける。\(t \in \mathbb{Z}\)
は離散流のパラメータであり、物理的時間ではない。

\subsubsection{§8.2 離散流}\label{ux96e2ux6563ux6d41}

自己更新写像
\(U : (\mathbf{k}, \Delta\mathbf{k}) \mapsto (\mathbf{k} + \Delta\mathbf{k}, \Delta\mathbf{k})\)
は \(\mathbb{Z}^4 \times \mathbb{Z}^4\)
上の離散流を生成する。波自身が「ラベルをどう進めるか」の規則
\(\Delta\mathbf{k}\) を内蔵し、外部の時間概念は流のパラメータ \(t\)
としてのみ現れる。

\subsubsection{§8.3
4次元固有の構造}\label{ux6b21ux5143ux56faux6709ux306eux69cbux9020}

拡張ラベル空間 \(\mathbb{Z}^4 \times \mathbb{Z}^4\) は Hurwitz
環の二重構造 \(\mathcal{H} \times \mathcal{H}\) を持つ。自己更新流が
\(\mathrm{SO}(4)\)
作用として整理され、\((\mathbf{k}, \Delta\mathbf{k})\)
の四元数積による演算が \(\mathbb{Z}^4 \times \mathbb{Z}^4\) 内に閉じる。

\begin{center}\rule{0.5\linewidth}{0.5pt}\end{center}

\subsection{§9 自己帰還ループと Nyquist
安定性}\label{ux81eaux5df1ux5e30ux9084ux30ebux30fcux30d7ux3068-nyquist-ux5b89ux5b9aux6027}

\subsubsection{§9.1
直交相互作用波の発出}\label{ux76f4ux4ea4ux76f8ux4e92ux4f5cux7528ux6ce2ux306eux767aux51fa}

波 \(\psi_{(\mathbf{k}, \Delta\mathbf{k})}\) の中心位相から、波数
\(\mathbf{k}\) に対する4次元直交方向に相互作用波 \(\chi\)
が発出される。\(\chi\) は (D1)--(D4) の球面波正則性により
\(\mathbb{Z}^4\) 上で形不変に広がる4次元球面波である。

\subsubsection{§9.2
自己再交差条件}\label{ux81eaux5df1ux518dux4ea4ux5deeux6761ux4ef6}

波の軌跡 \(\tau \mapsto \mathbf{x}_{\mathrm{self}}(\tau)\) と現在位置
\(\mathbf{x}_{\mathrm{self}}(t)\) について、再交差条件

\[|\mathbf{x}_{\mathrm{self}}(t) - \mathbf{x}_{\mathrm{self}}(\tau)| = c(t - \tau)\]

を満たす \(\tau < t\) が存在するとき、過去に発出された \(\chi_\tau\)
が現在の \(\psi\)
に乗法的に作用する。波が静止（\(\Delta\mathbf{k} = 0\)）なら \(\chi\)
は外向きに広がるだけで自身には戻らない。波が移動中（\(\Delta\mathbf{k} \neq 0\)）のときにのみ自己再交差が生じ、ラベル変換
\((\mathbf{k}, \Delta\mathbf{k}) \to (\mathbf{k}', \Delta\mathbf{k}')\)
が起こる。

\subsubsection{§9.3
ナイキスト安定性と速度上限}\label{ux30caux30a4ux30adux30b9ux30c8ux5b89ux5b9aux6027ux3068ux901fux5ea6ux4e0aux9650}

この自己帰還ループを制御工学的に分析する。開ループ伝達関数の位相構造:

\begin{itemize}
\tightlist
\item
  \textbf{基本位相}: \(-\pi/2\)（直交方向への発出）
\item
  \textbf{遅延位相}: \(-\omega \Delta t\)（再交差遅延）
\item
  \textbf{合成位相}: \(-\pi/2 - \omega \Delta t\)
\end{itemize}

4次元球面波の振幅減衰は \(1/r^{3/2}\)
であり、ループゲインが距離とともに十分速く減衰する。Nyquist 軌跡は
\((-1, 0)\) 点を囲まず、ループは安定する。

\textbf{速度上限の導出}: 有効遅延
\(\Delta t \propto 1/(c - v_\parallel)\) より、波の移動速度 \(v\)
が伝播速度 \(c\) に近づくと \(\Delta t \to 0\) となり、発散境界周波数
\(\omega_c = \pi/(2\Delta t) \to \infty\) となる。\(v = c\)
で全周波数が発散境界に乗る。

\textbf{命題 9.1}: Nyquist 安定性の維持条件として、\(v \leq c\)
が幾何学的に導出される。速度上限と変位上限は同じ安定境界
\(\angle G = -\pi, |G| = 1\)
への異なる接近経路として統一的に記述される。

\begin{center}\rule{0.5\linewidth}{0.5pt}\end{center}

\subsection{§10
中心投影と非線形効果の幾何学的起源}\label{ux4e2dux5fc3ux6295ux5f71ux3068ux975eux7ddaux5f62ux52b9ux679cux306eux5e7eux4f55ux5b66ux7684ux8d77ux6e90}

\subsubsection{§10.1
線形系の限界}\label{ux7ddaux5f62ux7cfbux306eux9650ux754c}

単一の主観空間 \(\mathbb{Z}^4\)
内では、形不変移動波は線形重ね合わせの閉性により
\textbf{相互作用なしで共存する}。加算操作のみではラベル集合の変換は起きない（Fourier
分解の一意性）。

\subsubsection{§10.2
乗法的作用の4つの方法}\label{ux4e57ux6cd5ux7684ux4f5cux7528ux306e4ux3064ux306eux65b9ux6cd5}

線形系に非線形効果を導入する方法:

\textbf{(i) 位相空間での二重性}: \(U(1)\)
値関数として系を定式化し、位相加算を振幅乗算として読む。

\textbf{(ii) 評価点格子の変換}: 評価点空間に四元数回転
\(\mathbf{x} \mapsto \mathbf{x} \cdot u\) を施し、ラベル空間に乗法的作用
\(\mathbf{k} \mapsto \mathbf{k} \cdot u^*\) を誘導する。\(n = 4\) 固有。

\textbf{(iii) 二系の位相合成}: 二つの独立な系の同点評価による位相合成。

\textbf{(iv) 親空間からの中心投影}: 親空間 \(S^5(R)\)
上の線形な波を、中心投影 \(\pi : S^5 \to \mathbb{Z}^4\)
で子空間に射影する。射影写像が分数写像（非線形）であるため、子空間上では
\textbf{非線形位相} を持つ関数として現れる。

\subsubsection{§10.3
中心投影の特権性}\label{ux4e2dux5fc3ux6295ux5f71ux306eux7279ux6a29ux6027}

\begin{enumerate}
\def\labelenumi{(\roman{enumi})}
\setcounter{enumi}{3}
\tightlist
\item
  は他の3方法と質的に異なる。(i)--(iii) は子空間 \(\mathbb{Z}^4\)
  の内部構造を利用し、子空間の対称性の範囲内で変化を起こす。(iv) は
  \textbf{枠の外} からの作用であり:
\end{enumerate}

\begin{itemize}
\tightlist
\item
  親空間での操作は単なる加算（線形）
\item
  子空間上では起源不明の非線形変化として現れる
\item
  子空間の対称性そのものを「外側から」変形できる
\end{itemize}

\textbf{命題 10.1}:
線形系に非線形効果を、力学を持ち込まずに幾何学的に正当化される形で導入する唯一の方法が、より高次元の親空間からの中心投影である。

\begin{center}\rule{0.5\linewidth}{0.5pt}\end{center}

\subsection{§11 de Sitter
時空への収束}\label{de-sitter-ux6642ux7a7aux3078ux306eux53ceux675f}

\subsubsection{§11.1
閉じた球面上のソリトン安定性}\label{ux9589ux3058ux305fux7403ux9762ux4e0aux306eux30bdux30eaux30c8ux30f3ux5b89ux5b9aux6027}

姉妹論文{[}1{]}で示した通り、閉じた \(n\) 次元球面 \(S^n\)
上で形不変定常波が安定に存在できるのは
\textbf{奇数次元に限られる}。偶数次元球面では毛玉定理（零点のないベクトル場が存在しない）により定常波構造が破綻する。

\subsubsection{§11.2
二つの議論の交差}\label{ux4e8cux3064ux306eux8b70ux8ad6ux306eux4ea4ux5dee}

{\def\LTcaptype{none} % do not increment counter
\begin{longtable}[]{@{}
  >{\raggedright\arraybackslash}p{(\linewidth - 6\tabcolsep) * \real{0.2000}}
  >{\raggedright\arraybackslash}p{(\linewidth - 6\tabcolsep) * \real{0.3000}}
  >{\raggedright\arraybackslash}p{(\linewidth - 6\tabcolsep) * \real{0.3000}}
  >{\raggedright\arraybackslash}p{(\linewidth - 6\tabcolsep) * \real{0.2000}}@{}}
\toprule\noalign{}
\begin{minipage}[b]{\linewidth}\raggedright
議論
\end{minipage} & \begin{minipage}[b]{\linewidth}\raggedright
対象空間
\end{minipage} & \begin{minipage}[b]{\linewidth}\raggedright
特権次元
\end{minipage} & \begin{minipage}[b]{\linewidth}\raggedright
根拠
\end{minipage} \\
\midrule\noalign{}
\endhead
\bottomrule\noalign{}
\endlastfoot
本論文（離散側） & 無限平坦格子 \(\mathbb{Z}^n\) & \(n = 4\) &
(F4)∩(F5)=\{4\} \\
姉妹論文（連続側） & 閉じた球面 \(S^n\) & 奇数次元 & Huygens 原理 +
毛玉定理 \\
\end{longtable}
}

両議論は独立に成立するが、物理的時空の構成に制約を課す:

\begin{itemize}
\tightlist
\item
  \textbf{空間部分}: 閉じた球面としてソリトンが安定 → \textbf{奇数次元}
  → 最小は \(S^3\)
\item
  \textbf{時間部分}: 無限に開いた方向で形不変波が伝播 →
  \textbf{双曲線的1次元}
\item
  \textbf{全体}: 局所平坦な4次元で (D1)--(D4) が成立 → \(3 + 1\) 次元
\end{itemize}

\subsubsection{§11.3 de Sitter
時空の幾何学的必然性}\label{de-sitter-ux6642ux7a7aux306eux5e7eux4f55ux5b66ux7684ux5fc5ux7136ux6027}

de Sitter 時空 \(\mathrm{dS}_4\) は5次元 Minkowski 空間
\(\mathbb{M}^{1,4}\) 内の双曲面

\[-X_0^2 + X_1^2 + X_2^2 + X_3^2 + X_4^2 = R^2\]

として定義される。空間部分は閉じた
\(S^3\)（奇数次元、ソリトン安定）、時間部分は双曲線（無限に開く、形不変波の伝播場）である。

\textbf{命題 11.1}: 形不変移動波の存在条件（離散側: \(n = 4\)
が唯一）と閉じたソリトンの存在条件（連続側:
奇数次元のみ）を同時に満たす4次元時空構造は、空間 \(S^3\)
と双曲線時間の積構造として記述される de Sitter 時空 \(\mathrm{dS}_4\)
のみである。

\subsubsection{§11.4
親子空間の役割分担}\label{ux89aaux5b50ux7a7aux9593ux306eux5f79ux5272ux5206ux62c5}

中心投影シリーズの枠組みにおいて:

\begin{itemize}
\tightlist
\item
  \textbf{子空間} \(S^4(R)\)（偶数次元、内部は \(\mathbb{Z}^4\)）:
  形不変移動波の運動場。表面では閉じたソリトン不安定。
\item
  \textbf{親空間} \(S^5(R)\)（奇数次元）:
  閉じたソリトンの供給源。中心投影で子空間に非線形効果を導入。
\end{itemize}

親空間が奇数次元、子空間の空間部分が奇数次元であるという構造は、奇偶次元の幾何学的差から
\textbf{必然的に} 決まる。

\begin{center}\rule{0.5\linewidth}{0.5pt}\end{center}

\subsection{§12 まとめ}\label{ux307eux3068ux3081}

本論文で示した論理の流れを整理する:

\begin{enumerate}
\def\labelenumi{\arabic{enumi}.}
\tightlist
\item
  \(\mathbb{Z}^n\) 上の球面定常波の存在条件 (D1)--(D4) は平方数次元
  \(\{4, 9, 16, \ldots\}\) で成立する（§3）。
\item
  形不変移動波の構造条件 (F1)--(F5)
  を加えると、\((F4) \cap (F5) = \{4\}\) により \(n = 4\)
  が唯一の次元として特定される（§4）。
\item
  \(\mathbb{Z}^4\) 上の形不変移動波は、保存量 (C1)--(C8)
  を持ち、粒子的条件 (G1)--(G5) を幾何学的に実現し、ラベル構造
  (L1)--(L5) で整理される（§5--§7）。
\item
  接束拡張 \(\mathbb{Z}^4 \times \mathbb{Z}^4\)
  により離散流が定義され、直交相互作用波の自己帰還ループが Nyquist
  安定性から速度上限 \(v \leq c\) を導出する（§8--§9）。
\item
  親空間からの中心投影が、線形系に非線形効果を時間概念なしで導入する唯一の幾何学的メカニズムである（§10）。
\item
  離散側（\(n = 4\)）と連続側（奇数次元球面）の二つの独立な議論が、de
  Sitter 時空 \(\mathrm{dS}_4\) を幾何学的必然として指し示す（§11）。
\end{enumerate}

本論文の主張はすべて純粋に幾何学的・数論的・代数的な事実として記述されており、物理的解釈を一切要求しない。物理的対応（粒子・相互作用・時空構造）は、これらの幾何学的事実に
\textbf{接続可能な解釈}
として開かれているが、命題本体の成立には影響しない。

\begin{center}\rule{0.5\linewidth}{0.5pt}\end{center}

\subsection{参考文献}\label{ux53c2ux8003ux6587ux732e}

{[}1{]} 木原範昭,
``閉じた球面上の形不変定常波------次元ごとの安定性と3+1時空の幾何学的必然性,''
2026.

{[}2{]} 木原範昭, ``中心投影による宇宙の3層モデル,'' Zenodo, 2025. DOI:
10.5281/zenodo.15083729.

{[}3{]} 木原範昭, ``球面間射影(グノモン正写像)の正則性,'' Zenodo, 2025.
DOI: 10.5281/zenodo.15292757.

{[}4{]} 木原範昭, ``(R,Q)マッピング定義論文{[}M3{]},'' Zenodo, 2025.
DOI: 10.5281/zenodo.19692853.
\end{document}
